\documentclass[10pt,a4paper]{amsart}
\usepackage{amsmath,amssymb,amsthm,graphicx}
% Retained minimal preamble for the Stacks Project chapter flat.tex (More on Flatness).
% Verbatim excerpt of the chapter preamble: Xy-pic setup (native preamble lines 14-18) and the native
% theorem declarations with the mathematical macros (native preamble lines 154-217). Copied mechanically;
% nothing added, removed or redefined. The wrapper class and the amsmath/amssymb/amsthm/graphicx base
% packages are supplied by the assembler, so the class/externaldocument part of the native preamble is omitted.
% For commutative diagrams we use Xy-pic
\usepackage[all]{xy}

% We use 2cell for 2-commutative diagrams.
\xyoption{2cell}
% Theorem environments.
%
\theoremstyle{plain}
\newtheorem{theorem}[subsection]{Theorem}
\newtheorem{proposition}[subsection]{Proposition}
\newtheorem{lemma}[subsection]{Lemma}

\theoremstyle{definition}
\newtheorem{definition}[subsection]{Definition}
\newtheorem{example}[subsection]{Example}
\newtheorem{exercise}[subsection]{Exercise}
\newtheorem{situation}[subsection]{Situation}

\theoremstyle{remark}
\newtheorem{remark}[subsection]{Remark}
\newtheorem{remarks}[subsection]{Remarks}

\numberwithin{equation}{subsection}

% Macros
%
\def\lim{\mathop{\mathrm{lim}}\nolimits}
\def\colim{\mathop{\mathrm{colim}}\nolimits}
\def\Spec{\mathop{\mathrm{Spec}}}
\def\Hom{\mathop{\mathrm{Hom}}\nolimits}
\def\Ext{\mathop{\mathrm{Ext}}\nolimits}
\def\SheafHom{\mathop{\mathcal{H}\!\mathit{om}}\nolimits}
\def\SheafExt{\mathop{\mathcal{E}\!\mathit{xt}}\nolimits}
\def\Sch{\mathit{Sch}}
\def\Mor{\mathop{\mathrm{Mor}}\nolimits}
\def\Ob{\mathop{\mathrm{Ob}}\nolimits}
\def\Sh{\mathop{\mathit{Sh}}\nolimits}
\def\NL{\mathop{N\!L}\nolimits}
\def\CH{\mathop{\mathrm{CH}}\nolimits}
\def\proetale{{pro\text{-}\acute{e}tale}}
\def\etale{{\acute{e}tale}}
\def\QCoh{\mathit{QCoh}}
\def\Ker{\mathop{\mathrm{Ker}}}
\def\Im{\mathop{\mathrm{Im}}}
\def\Coker{\mathop{\mathrm{Coker}}}
\def\Coim{\mathop{\mathrm{Coim}}}

% Boxtimes
%
\DeclareMathSymbol{\boxtimes}{\mathbin}{AMSa}{"02}

%
% Macros for moduli stacks/spaces
%
\def\QCohstack{\mathcal{QC}\!\mathit{oh}}
\def\Cohstack{\mathcal{C}\!\mathit{oh}}
\def\Spacesstack{\mathcal{S}\!\mathit{paces}}
\def\Quotfunctor{\mathrm{Quot}}
\def\Hilbfunctor{\mathrm{Hilb}}
\def\Curvesstack{\mathcal{C}\!\mathit{urves}}
\def\Polarizedstack{\mathcal{P}\!\mathit{olarized}}
\def\Complexesstack{\mathcal{C}\!\mathit{omplexes}}
% \Pic is the operator that assigns to X its picard group, usage \Pic(X)
% \Picardstack_{X/B} denotes the Picard stack of X over B
% \Picardfunctor_{X/B} denotes the Picard functor of X over B
\def\Pic{\mathop{\mathrm{Pic}}\nolimits}
\def\Picardstack{\mathcal{P}\!\mathit{ic}}
\def\Picardfunctor{\mathrm{Pic}}
\def\Deformationcategory{\mathcal{D}\!\mathit{ef}}

% --- appended 2026-10-12 (Chief): the project own definitions of the reference/slogan/history
% display environments, copied mechanically from the shared Stacks preamble lines 8-12
% (\usepackage{verbatim} plus the three \newenvironment lines, which the chapter preamble keeps
% verbatim from the project preamble). Nothing added, removed or redefined: these environments
% carry no mathematical content and the project itself typesets them as comments. Without these
% three lines a counted statement whose own text uses them fails with Environment slogan undefined.
\usepackage{verbatim}
\newenvironment{reference}{\comment}{\endcomment}
\newenvironment{slogan}{\comment}{\endcomment}
\newenvironment{history}{\comment}{\endcomment}
\usepackage{cleveref}
\crefname{originallemma}{original lemma}{original lemmas}
\Crefname{originallemma}{Original lemma}{Original lemmas}
\makeatletter
\newlabel{lemma-colim-F-Vect}{{0EXB}{source}{}{sourcepointer-1}{}}
\newlabel{lemma-colim-F-Vect@cref}{{[originallemma][0][]0EXB}{[1][0][]source}{}{}{}}
\makeatother
\crefname{originallemma}{original lemma}{original lemmas}
\Crefname{originallemma}{Original lemma}{Original lemmas}
\makeatletter
\newlabel{lemma-refine-check-h-stack}{{0EX9}{source}{}{sourcepointer-2}{}}
\newlabel{lemma-refine-check-h-stack@cref}{{[originallemma][0][]0EX9}{[1][0][]source}{}{}{}}
\makeatother
\crefname{originalproposition}{original proposition}{original propositions}
\Crefname{originalproposition}{Original proposition}{Original propositions}
\makeatletter
\newlabel{descent-proposition-fpqc-descent-quasi-coherent}{{023T}{source}{}{sourcepointer-3}{}}
\newlabel{descent-proposition-fpqc-descent-quasi-coherent@cref}{{[originalproposition][0][]023T}{[1][0][]source}{}{}{}}
\makeatother
\crefname{originallemma}{original lemma}{original lemmas}
\Crefname{originallemma}{Original lemma}{Original lemmas}
\makeatletter
\newlabel{descent-lemma-finite-locally-free-descends}{{05B2}{source}{}{sourcepointer-4}{}}
\newlabel{descent-lemma-finite-locally-free-descends@cref}{{[originallemma][0][]05B2}{[1][0][]source}{}{}{}}
\makeatother
\crefname{originallemma}{original lemma}{original lemmas}
\Crefname{originallemma}{Original lemma}{Original lemmas}
\makeatletter
\newlabel{lemma-vector-bundle-I}{{0EXC}{source}{}{sourcepointer-5}{}}
\newlabel{lemma-vector-bundle-I@cref}{{[originallemma][0][]0EXC}{[1][0][]source}{}{}{}}
\makeatother
\crefname{originallemma}{original lemma}{original lemmas}
\Crefname{originallemma}{Original lemma}{Original lemmas}
\makeatletter
\newlabel{lemma-vector-bundle-II}{{0EXD}{source}{}{sourcepointer-6}{}}
\newlabel{lemma-vector-bundle-II@cref}{{[originallemma][0][]0EXD}{[1][0][]source}{}{}{}}
\makeatother
\title{Vector bundles for the colimit along Frobenius form a stack in groupoids for the h topology}
\author{The Stacks Project Authors}
\date{stacks-project:flat.tex:local-snapshot-20261011; GNU FDL 1.2 or later}
\begin{document}
\maketitle
\noindent\textit{Editorial note.} All delivered mathematical text is the author's own, copied from the pinned Stacks Project chapter file flat.tex; no statement or proof step was rewritten, reordered or removed.\par\medskip
\noindent\textit{Editorial note.} The preamble is the chapter's own retained minimal preamble (sources/stacks-flat/retained-preamble.tex), a verbatim excerpt of the native Xy-pic setup and of the native theorem declarations and mathematical macros, followed by the project own definitions of the reference, slogan and history display environments (verbatim package plus three newenvironment lines) that the retained file carries at its end and that the project itself typesets as comments. No macro was invented or redefined, and the wrapper class and the amsmath/amssymb/amsthm/graphicx base packages are supplied by the assembler.\par\medskip
\noindent\textit{Editorial note.} The counted statement and proof contain no cross reference, no citation and no included file, so the delivered body is self contained.\par\medskip
\noindent\textit{Editorial note.} Canonical location of the counted proposition: Stacks tag 0EXE (sources/stacks-flat/pinned-original-source.tex line 12440).\par\medskip
\noindent\textit{Allowed standard prerequisites (editorial).}\begin{itemize}
\item The stack conditions to be checked. Same chapter, Lemma labelled lemma-refine-check-h-stack, cited at proof lines 12456-12458. Lists conditions (1), (2) and (3).
\item fppf descent of finite locally free modules. Descent, Proposition labelled descent-proposition-fpqc-descent-quasi-coherent and Lemma labelled descent-lemma-finite-locally-free-descends, cited at proof lines 12468-12472. Supplies the equivalence of categories verifying condition (2).
\item The fibre category as a colimit of vector bundle categories. Same chapter, Lemma labelled lemma-colim-F-Vect, cited at proof lines 12452-12454. Gives the description of the fibre categories used throughout.
\item The two almost blowup computations. Same chapter, Lemmas labelled lemma-vector-bundle-I and lemma-vector-bundle-II, cited at proof lines 12478-12483. Settle condition (3).
\end{itemize}\medskip
\ifdefined\hypertarget\hypertarget{sourcepointer-1}{}\fi
\noindent\textit{Source reference.} Original lemma 0EXB: sources/stacks-flat/pinned-original-source.tex, source lines 12223--12223. Same-chapter pointer into the author environment lemma at source line 12223; invoked at proof line 12453 of the counted proof. It is retained only as a navigation pointer; every load-bearing step is inside the delivered span.\par\medskip
\ifdefined\hypertarget\hypertarget{sourcepointer-2}{}\fi
\noindent\textit{Source reference.} Original lemma 0EX9: sources/stacks-flat/pinned-original-source.tex, source lines 11696--11696. Same-chapter pointer into the author environment lemma at source line 11696; invoked at proof line 12455 of the counted proof. It is retained only as a navigation pointer; every load-bearing step is inside the delivered span.\par\medskip
\ifdefined\hypertarget\hypertarget{sourcepointer-3}{}\fi
\noindent\textit{Source reference.} Original proposition 023T: sources/stacks-chow/tags.txt, source lines 2754--2754. Cross-chapter standard prerequisite located in the project tag list staged in the local snapshot; invoked at proof line 12477 of the counted proof. It is retained only as a navigation pointer; every load-bearing step is inside the delivered span.\par\medskip
\ifdefined\hypertarget\hypertarget{sourcepointer-4}{}\fi
\noindent\textit{Source reference.} Original lemma 05B2: sources/stacks-chow/tags.txt, source lines 6870--6870. Cross-chapter standard prerequisite located in the project tag list staged in the local snapshot; invoked at proof line 12478 of the counted proof. It is retained only as a navigation pointer; every load-bearing step is inside the delivered span.\par\medskip
\ifdefined\hypertarget\hypertarget{sourcepointer-5}{}\fi
\noindent\textit{Source reference.} Original lemma 0EXC: sources/stacks-flat/pinned-original-source.tex, source lines 12235--12235. Same-chapter pointer into the author environment lemma at source line 12235; invoked at proof line 12482 of the counted proof. It is retained only as a navigation pointer; every load-bearing step is inside the delivered span.\par\medskip
\ifdefined\hypertarget\hypertarget{sourcepointer-6}{}\fi
\noindent\textit{Source reference.} Original lemma 0EXD: sources/stacks-flat/pinned-original-source.tex, source lines 12303--12303. Same-chapter pointer into the author environment lemma at source line 12303; invoked at proof line 12482 of the counted proof. It is retained only as a navigation pointer; every load-bearing step is inside the delivered span.\par\medskip
\begin{proposition}
\label{proposition-h-descent-vector-bundles-p}
Let $p$ be a prime number. Let $S$ be a scheme in characteristic $p$.
Then the category fibred in groupoids
$$
p : \mathcal{S} \longrightarrow (\Sch/S)_h
$$
whose fibre category over $U$ is the category
of finite locally free $\colim_F \mathcal{O}_U$-modules over $U$
is a stack in groupoids. Moreover, if $U$ is quasi-compact
and quasi-separated, then $\mathcal{S}_U$ is $\colim_F \textit{Vect}(U)$.
\end{proposition}

\begin{proof}
The final assertion is the content of Lemma \ref{lemma-colim-F-Vect}.
To prove the proposition we will check conditions (1), (2), and (3) of
Lemma \ref{lemma-refine-check-h-stack}.

\medskip\noindent
Condition (1) holds because by definition we have glueing
for the Zariski topology.

\medskip\noindent
To see condition (2), suppose that $f : X \to Y$ is a surjective,
flat, proper morphism of finite presentation over $S$ with $Y$ affine.
Since $Y, X, X \times_Y X$ are quasi-compact and quasi-separated,
we can use the description of fibre categories given in the
statement of the proposition. Then it is clearly enough to
show that
$$
\textit{Vect}(Y)
\longrightarrow
\textit{Vect}(X)
\times_{\textit{Vect}(X \times_Y X)}
\textit{Vect}(X)
$$
is an equivalence (as this will imply the same for the colimits).
This follows immediately from fppf descent of finite locally free modules, see
Descent, Proposition \ref{descent-proposition-fpqc-descent-quasi-coherent} and
Lemma \ref{descent-lemma-finite-locally-free-descends}.

\medskip\noindent
Condition (3) is the content of
Lemmas \ref{lemma-vector-bundle-I} and \ref{lemma-vector-bundle-II}.
\end{proof}

\end{document}
